% Appendix A — Data (first appendix section).
\subsection{MSSV cohort and quality screening}
\label{app:data:cohort}

MSSV (OpenNeuro ds006366)~\cite{rose2025mssv} comprises invasive EEG and neck EMG from 92 mice across five European laboratories, with expert Wake, NREM, REM, and Artifact labels on a 4\,s grid for evaluation only.
This report analyses laboratories 2 (Bern), 3 (Copenhagen), and 5 (Lyon): healthy wild-type sites with different electrode layouts, excluding laboratories 1 and 4 (pathology-specific or single-electrode cohorts).

Quality screening retains mice with low EEG drift and EMG that separates Wake from NREM.
Drift is quantified as relative 0.5--2\,Hz power $\leq 0.2$ of 0.5--30\,Hz on each retained electrode.
Wake/NREM separability uses a root-mean-square ratio $\geq 2$ on neck EMG.
The screened cohort comprises 20 mice (6 Bern, 10 Copenhagen, 4 Lyon) with 42 recording runs and ${\sim}453$\,h of EEG/EMG (${\sim}4.1 \times 10^{5}$ artefact-free 4\,s epochs after screening).
Training mice never appear in the validation fold for a given split (\apdx{app:data:splits}).
No new animal procedures were performed.

\apptabdesc{Quality-screened cohort inventory. Per-mouse recording runs, hours, and 4\,s epoch counts (MSSV annotation grid).}
\label{app:tab:cohort}
\apptabular{\small% Auto-generated by scripts/paper/build_cv4fold_split_tables.py
\begin{tabular}{|l|l|r|r|r|}
\hline
\textbf{Mouse} & \textbf{Site} & \textbf{Runs} & \textbf{Hours} & \textbf{Epochs (4\,s)} \\
\thickhline
sub-071 & Bern (2) & 2 & 12 & 10{,}800 \\
sub-072 & Bern (2) & 2 & 12 & 10{,}800 \\
sub-076 & Bern (2) & 2 & 12 & 10{,}800 \\
sub-077 & Bern (2) & 2 & 13 & 11{,}700 \\
sub-080 & Bern (2) & 2 & 12 & 10{,}800 \\
sub-081 & Bern (2) & 2 & 12 & 10{,}800 \\
sub-038 & Copenhagen (3) & 2 & 48 & 43{,}201 \\
sub-039 & Copenhagen (3) & 3 & 12 & 10{,}836 \\
sub-041 & Copenhagen (3) & 3 & 12 & 10{,}803 \\
sub-043 & Copenhagen (3) & 1 & 24 & 21{,}600 \\
sub-048 & Copenhagen (3) & 3 & 12 & 10{,}804 \\
sub-054 & Copenhagen (3) & 2 & 48 & 43{,}194 \\
sub-056 & Copenhagen (3) & 1 & 23.6 & 21{,}196 \\
sub-059 & Copenhagen (3) & 1 & 24 & 21{,}600 \\
sub-060 & Copenhagen (3) & 1 & 8 & 7{,}212 \\
sub-069 & Copenhagen (3) & 1 & 24 & 21{,}600 \\
sub-087 & Lyon (5) & 6 & 72 & 64{,}788 \\
sub-088 & Lyon (5) & 2 & 24 & 21{,}595 \\
sub-089 & Lyon (5) & 2 & 24 & 21{,}596 \\
sub-092 & Lyon (5) & 2 & 24 & 21{,}595 \\
\hline
\textbf{Total (20 mice)} &  & 42 & 452.6 & 407{,}320 \\
\hline
\end{tabular}
}

\subsection{Preprocessing and spectral inputs}
\label{app:data:preprocessing}

Montages are harmonised to a parietal-frontal EEG differential plus neck EMG.
Bern (laboratory~2) provides four EEG contacts (EEG1/2 parietal, EEG3/4 frontal); we use EEG1 and EEG4 to match the single parietal--frontal pairing at laboratories~3 and~5 (EEG1 parietal, EEG2 frontal).
Recordings are segmented into 4\,s epochs at 128\,Hz. Artifact epochs are removed.
Each epoch is a log-power spectral map $\mathbf{x}_t \in \mathbb{R}^{C \times F}$, batched into contiguous sequences for temporal models without hand-crafted expert features.

Per-laboratory channel sets were fixed from incohort ablation winners (full sweep outcomes in \apdx{app:incohort_ceiling}).
Spectral passbands are summarised in \apdx{app:data:bands}.

\subsection{Spectral bands and frequency content}
\label{app:data:bands}

Each 4\,s epoch comprises 512 samples at 128\,Hz sampling rate.
The CNN input is a log-power spectrum from an unwindowed real FFT (frequency resolution 0.25\,Hz; Nyquist 64\,Hz).
Band limiting is applied in the frequency domain on the complex spectrum before log-power: a null lower cutoff retains all bins from DC upward, and an upper cutoff zeroes bins above the configured limit.
The model therefore receives the full spectral shape inside each passband, not a small vector of hand-defined sleep bands.

\paragraph*{Locked per-laboratory passbands (holdout protocol).}
\apptabular{\small
\begin{tabular}{|l|l|l|l|}
\hline
\textbf{Site} & \textbf{EEG channels} & \textbf{EEG passband} & \textbf{EMG passband} \\
\thickhline
Bern (lab~2) & EEG1, EEG4 & 0--30\,Hz & 3--100\,Hz ($\leq 64$\,Hz Nyquist) \\
Copenhagen (lab~3) & EEG1, EEG2 & 0--20\,Hz & 5--60\,Hz \\
Lyon (lab~5) & EEG1, EEG2 & 0--20\,Hz & 5--60\,Hz \\
\hline
\end{tabular}}

\paragraph*{Diagnostic and physiology bands.}
Quality screening compares relative 0.5--2\,Hz EEG power to 0.5--30\,Hz total (\apdx{app:data:cohort}).
For substage physiology panels (Fig~\ref{fig3} and Fig~\ref{fig4}) and summary tables, log-power spectra are aggregated into interpretable bands on the same frequency grid:
delta 0.5--4\,Hz, theta 4--8\,Hz, and broadband EMG power summed over the retained EMG passband.
Reported $\theta/\delta$ ratios use these bands.

Bern uses no post-normalisation on log-power spectra.
Copenhagen and Lyon apply run-wise z-scoring after band limiting: for each recording run, each channel's log-power spectrum is scaled by the mean and standard deviation computed across all 4\,s epochs in that run (post-normalisation).
These limits were chosen from incohort ablations (\apdx{app:incohort_ceiling}) and held fixed for all cross-laboratory holdout runs.

\subsection{Cross-validation protocol}
\label{app:data:splits}

Folds were designed so each site contributes holdout mice every fold while balancing validation load: fold~1 is epoch-heaviest (${\sim}32\%$ of cohort epochs in validation, driven by long Copenhagen and Lyon recordings) whereas folds~2--4 hold out 21--25\% of epochs each.
We use four folds because Lyon contributes only four quality-screened mice, which is the maximum fold count for strict mouse-level holdout.
sub-069 is assigned to fold~2 rather than fold~4 so Copenhagen does not contribute four simultaneous holdout mice on the substage-analysis fold.
Each VAE configuration trains three random seeds per fold. Reported metrics use the best prior NMI among seeds.

\apptabdesc{Four-fold leave-mice-out assignments. Top: held-out mice per site and fold.
Bottom: validation versus training pool sizes (joint scope trains on all non-held-out mice across sites).}
\label{app:tab:splits}
\apptabular{% Auto-generated by scripts/paper/build_cv4fold_split_tables.py
\begin{tabular}{|c|l|l|r|r|r|}
\hline
\textbf{Fold} & \textbf{Site} & \textbf{Held-out mice} & \textbf{Runs} & \textbf{Hours} & \textbf{Epochs} \\
\thickhline
1 & Bern (2) & sub-071 & 2 & 12 & 10{,}800 \\
 & Copenhagen (3) & sub-038, sub-039 & 5 & 60 & 54{,}037 \\
 & Lyon (5) & sub-087 & 6 & 72 & 64{,}788 \\
\hline
2 & Bern (2) & sub-072, sub-076 & 4 & 24 & 21{,}600 \\
 & Copenhagen (3) & sub-041, sub-048, sub-069 & 7 & 48 & 43{,}207 \\
 & Lyon (5) & sub-088 & 2 & 24 & 21{,}595 \\
\hline
3 & Bern (2) & sub-077 & 2 & 13 & 11{,}700 \\
 & Copenhagen (3) & sub-043, sub-054 & 3 & 72 & 64{,}794 \\
 & Lyon (5) & sub-089 & 2 & 24 & 21{,}596 \\
\hline
4 & Bern (2) & sub-080, sub-081 & 4 & 24 & 21{,}600 \\
 & Copenhagen (3) & sub-056, sub-059, sub-060 & 3 & 55.6 & 50{,}008 \\
 & Lyon (5) & sub-092 & 2 & 24 & 21{,}595 \\
\hline
\end{tabular}
\par\medskip
\begin{tabular}{|c|r|r|r|r|r|r|r|}
\hline
\textbf{Fold} & \textbf{Val mice} & \textbf{Val runs} & \textbf{Val hours} & \textbf{Val epochs} & \textbf{Train mice} & \textbf{Train epochs} & \textbf{Val (\%)} \\
\thickhline
1 & 4 & 13 & 144 & 129{,}625 & 16 & 277{,}695 & 31.8 \\
2 & 6 & 13 & 96 & 86{,}402 & 14 & 320{,}918 & 21.2 \\
3 & 4 & 7 & 109 & 98{,}090 & 16 & 309{,}230 & 24.1 \\
4 & 6 & 9 & 103.5 & 93{,}203 & 14 & 314{,}117 & 22.9 \\
\hline
\end{tabular}
}
